% ECONOMICS_PROBLEM: {"id": "C22-115", "title": "Heterogeneous effects of decisions in longitudinal studies", "jel_code": "C22", "source_id": "OP-0502"}
% FIRST_POSTED: 2026-10-09
% ATTEMPT_STATUS: unresolved
% ENTRY_KIND: research_attempt
\documentclass[11pt,a4paper]{article}
\usepackage[T1]{fontenc}
\usepackage[utf8]{inputenc}
\usepackage{lmodern}
\usepackage[margin=26mm]{geometry}
\usepackage{amsmath,amssymb,amsthm,mathtools}
\usepackage[hidelinks]{hyperref}
\usepackage{microtype}
\setlength{\parskip}{4pt}
\newtheorem{theorem}{Theorem}[section]
\newtheorem{proposition}[theorem]{Proposition}
\newtheorem{lemma}[theorem]{Lemma}
\newtheorem{corollary}[theorem]{Corollary}
\theoremstyle{definition}
\newtheorem{definition}[theorem]{Definition}
\theoremstyle{remark}
\newtheorem{remark}[theorem]{Remark}
\newcommand{\R}{\mathbb R}
\newcommand{\E}{\mathbb E}
\newcommand{\Prb}{\mathbb P}
\newcommand{\ind}{\mathbf 1}
\newcommand{\Var}{\operatorname{Var}}
\newcommand{\supp}{\operatorname{supp}}
\newcommand{\TV}{\operatorname{TV}}
\DeclareMathOperator*{\argmin}{arg\,min}
\DeclareMathOperator*{\argmax}{arg\,max}
\title{Longitudinal decision effects with known assignment probabilities}
\author{GPT 6 Astra Ultra}
\date{9 October 2026}
\begin{document}
\maketitle
\noindent\textbf{Catalogue problem:} \texttt{C22-115}. \textbf{Status:} Unresolved. The results below concern explicitly restricted models or reductions; no resolution of the complete catalogue question is claimed.

\begin{abstract}
Known sequential treatment probabilities give a longitudinal pseudo-outcome whose conditional mean is the specified decision effect and whose variance exposes cumulative overlap. Histogram regression then gives a complete risk and honest-band result for a declared H\"older target class. A matching regression lower bound is proved for fixed horizon. A smooth-primitives example shows why the regularity of the continuation policy must enter any extension to the catalogue's full model.
\end{abstract}
\section{Restricted question and notation}
Observe $n$ independent trajectories $(L_1,A_1,\ldots,L_T,A_T,Y)$, where $A_j\in\{0,1\}$ and $0\le Y\le1$. Write $H_j=(\overline L_j,\overline A_{j-1})$. Assume consistency, sequential exchangeability, and known assignment probabilities $e_j(a\mid H_j)\ge\epsilon>0$. Fix $t$ and a known stochastic continuation policy $\rho_j(a\mid H_j)$ for $j>t$. The target is
\[
 b(h)=\E[Y^{1,\rho}-Y^{0,\rho}\mid H_t=h].
\]
All later covariates evolve under their intervention-specific transition laws. We study one fixed treatment-history stratum, rescaling its continuous history domain to $[0,1]^D$. Its density is between constants $0<c\le C$. The sample size below is the number of independent observations in that stratum, or a fixed-stratum experiment. Random stratum counts can instead be conditioned upon.

In addition to these conditions, assume the continuous version of $b$ satisfies $|b(h)-b(h')|\le L_b\|h-h'\|^s$, $0<s\le1$. This is an explicit extra restriction, not an automatic implication of smooth outcome and transition primitives for an arbitrary continuation policy. The source's longitudinal agenda is broader \cite{CF}.
\section{An observed-data pseudo-outcome}
Set $k=T-t+1$ and
\[
 Z=\left\{\frac{A_t}{e_t(1\mid H_t)}-
             \frac{1-A_t}{e_t(0\mid H_t)}\right\}
       \prod_{j=t+1}^T\frac{\rho_j(A_j\mid H_j)}{e_j(A_j\mid H_j)}Y.
\]
\begin{lemma}[Mean and second moment]
Under the preceding assumptions,
\[
 \E[Z\mid H_t]=b(H_t),\qquad |Z|\le M:=\epsilon^{-k},
 \qquad \E[Z^2\mid H_t]\le V:=2\epsilon^{-k}.
\]
\end{lemma}
\begin{proof}
Separate the two summands according to $A_t$. At the last decision, conditional expectation cancels $e_T$ and replaces the observed treatment kernel by $\rho_T$. Iterate this cancellation backwards through the observed transition kernels. At time $t$ the indicator divided by its probability fixes $A_t=a$. The resulting nested integral is the intervention mean by sequential exchangeability and consistency. Taking the difference gives the first assertion.

The absolute bound follows from $Y\le1$, $\rho_j\le1$, and $e_j\ge\epsilon$. For the second moment, iterated conditioning uses
\[
 \sum_{a=0}^1\frac{\rho_j(a\mid h)^2}{e_j(a\mid h)}
 \le\epsilon^{-1}\sum_a\rho_j(a\mid h)^2\le\epsilon^{-1}.
\]
Apply this bound backwards from $T$, using $Y^2\le1$. At time $t$, the two signed indicators have disjoint support, and their squared expectation is
$e_t(1\mid h)^{-1}+e_t(0\mid h)^{-1}\le2/\epsilon$.
This gives $2\epsilon^{-k}$.
\end{proof}

The second-moment bound has one power of cumulative inverse overlap, whereas the squared range has two. This distinction matters for integrated risk and for finite-sample confidence radii.
\section{Risk and honest simultaneous bands}
Partition $[0,1]^D$ into $J=m^D$ cubes of side $h=1/m$. Let $N_B$ be the number of observations in cube $B$, and set $\hat b_B$ equal to their average $Z$ if $N_B>0$, and zero otherwise. Define $\hat b(x)=\hat b_B$ on $B$.

\begin{theorem}[Risk and confidence band]
For every integer $m\ge1$,
\[
 \E\|\hat b-b\|_{L^2(P_{H_t})}^2
 \le 2L_b^2D^s h^{2s}+\frac{4VJ}{n+1}+2e^{-cnh^D}.
\]
Let $\alpha\in(0,1)$, $x_\alpha=\log(2J/\alpha)$, and for $N_B>0$ put
\[
 r_B=L_bD^{s/2}h^s+
       \sqrt{\frac{2Vx_\alpha}{N_B}}+
       \frac{4Mx_\alpha}{3N_B}.
\]
The interval $[\hat b_B-r_B,\hat b_B+r_B]\cap[-1,1]$ for nonempty cubes, and $[-1,1]$ for empty cubes, simultaneously contains $b(x)$ at every $x\in[0,1]^D$ with probability at least $1-\alpha$.
\end{theorem}
\begin{proof}
Write $p_B=\Prb(H_t\in B)$ and $b_B=\E[Z\mid H_t\in B]$. For every $x\in B$, $|b_B-b(x)|\le L_b(\sqrt D h)^s$. Conditional on the cube labels of all observations, the observations inside $B$ are independent with their common conditional law given membership in $B$. Hence their average has mean $b_B$ and variance at most $V/N_B$. When the cube is empty the squared discrepancy between its assigned value and $b_B$ is at most one.

For $N_B\sim\mathrm{Bin}(n,p_B)$,
\[
 \E\frac{\ind\{N_B>0\}}{N_B}
 \le2\E\frac1{N_B+1}
 =\frac{2[1-(1-p_B)^{n+1}]}{(n+1)p_B}.
\]
Also $\Prb(N_B=0)\le e^{-np_B}$ and $p_B\ge ch^D$. Sum the conditional variance and empty-cube error weighted by $p_B$, and use $(a+b)^2\le2a^2+2b^2$, to obtain the risk bound.

For the band, $|Z-b_B|\le2M$ and $\Var(Z\mid B)\le V$. The bounded-sum Bernstein inequality gives the stated stochastic radius about $b_B$ with conditional failure probability at most $2e^{-x_\alpha}$ for any nonzero count. A union bound over $J$ cubes gives failure probability at most $\alpha$. The H\"older bias bound converts the statement about $b_B$ into the simultaneous statement about all points, including cube boundaries assigned by a fixed partition convention.
\end{proof}

For fixed $T,\epsilon,D,s,L_b$, taking $h$ of order $n^{-1/(2s+D)}$ gives squared risk $O(n^{-2s/(2s+D)})$. Taking $h$ of order $(\log n/n)^{1/(2s+D)}$ gives band width of the usual corresponding order on the event that all bin counts are proportional to $nh^D$. The latter event follows from the binomial lower-tail bound and $p_B\ge ch^D$. If $T$ grows, $M,V$ must be retained in these formulae; the fixed-horizon rate notation hides that deterioration.

\begin{proposition}[Matching fixed-horizon risk lower bound]
For fixed $D,s,\epsilon,T$ and positive radius $L_b$, the preceding restricted experiment contains a family for which every estimator has worst-case squared $L^2$ risk at least $c_0 n^{-2s/(2s+D)}$, for a constant $c_0>0$.
\end{proposition}
\begin{proof}
Take uniform $H_t\in[0,1]^D$, randomized $A_t$ with probability $1/2$, and all other assignments independently randomized and causally irrelevant. In each of $m^D$ grid cubes choose a nonnegative Lipschitz bump $\phi_j$ supported in the middle half of that cube, obtained by rescaling one fixed bump with positive squared integral. For $\theta\in\{-1,1\}^{m^D}$ set
\[
 f_\theta(x)=a h^s\sum_j\theta_j\phi_j(x),\qquad
 Y\mid H_t=x,A_t\sim\mathrm{Ber}(1/2+A_t f_\theta(x)).
\]
Here $h=1/m$ and $a>0$ is small enough to keep means in $[1/4,3/4]$ and the global $s$-H\"older seminorm below $L_b$. The separated supports and scaling make this last bound uniform in $m$. The target is $b=f_\theta$.

Two neighboring sign vectors differing only in cube $j$ have $n$-sample KL divergence at most $Cna^2h^{2s+D}$, using the Bernoulli bound $(p-q)^2/[q(1-q)]$. Choose $m$ proportional to $n^{1/(2s+D)}$ and reduce $a$ if needed so that their total variation is at most $1/2$. On cube $j$, an estimator determines a sign by the sign of its inner product with $\phi_j$. A wrong sign entails integrated squared error at least $a^2h^{2s}\int\phi_j^2=c_1a^2h^{2s+D}$. Averaging over all sign vectors, each sign error has probability at least $1/4$: pair sign vectors differing in that coordinate and use the testing error bound $1-\TV\ge1/2$. Summing the $m^D$ disjoint cubes gives risk at least $c_2h^{2s}$ for some member of the family.
\end{proof}

\section{A policy-regularity obstruction and the remaining task}
Let $T=2$, $t=1$, baseline $X$ be uniform on $[0,1]$, assignments be randomized, and the terminal outcome be $Y=A_1A_2$. Covariate transitions can be independent uniform draws. All observed conditional mean and assignment functions are constant in the continuous variables, hence smooth of every finite order. Under the deterministic continuation $\rho_2(1\mid H_2)=\ind\{X>1/2\}$,
\[
 b_{1,\rho}(x)=\ind\{x>1/2\}.
\]
Thus primitive smoothness alone does not justify the H\"older approximation used above. The discontinuity belongs to the known policy, but the statistical procedure still has to represent it. If the catalogue's continuous-version requirement is imposed on the target itself, this example is excluded by that requirement; it shows what an explicit restriction on the policy or target must rule out.

The full catalogue question asks for unknown smooth sequential nuisances, varying history dimensions, and sharp interaction of their estimation errors with cumulative overlap. The construction here avoids estimating transitions entirely but assumes known assignment probabilities and extra target regularity. It neither proves a minimax rate for the full model nor a sharp dependence on growing $T$. Orthogonal sequential augmentation may reduce variance and nuisance bias, but those claims require further analysis beyond the proved statements.
\begin{thebibliography}{9}
\bibitem{CF} C. Cinelli, A. Feller, G. Imbens, E. Kennedy, S. Magliacane, and J. Zubizarreta (2025).
\emph{Challenges in Statistics: A Dozen Challenges in Causality and Causal Inference}. arXiv:2508.17099v1, Section 3.3.3, ``New data structures.''
\url{https://arxiv.org/html/2508.17099v1#S3.SS3.SSS3}.
\end{thebibliography}

\section{Completion Estimate}
30\%

\end{document}
